% kotexgweb.tex -- Korean localization and a LuaTeX PDF back end for GWEB.
%
% Put  \input kotexgweb.tex  in the limbo of a .w file and process the woven
% output with luatex (not pdftex):  luatex foo.tex.  gweave emits
% "\input gwebmac" as the very first line, so this file is read afterwards and
% overrides gwebmac wherever Korean text or the LuaTeX engine needs it.
%
% Requires luatexko (loaded through ko-font.tex) and the Noto Serif/Sans CJK KR
% fonts.  Korean font choices live in the companion file ko-font.tex; edit that
% file to change typefaces.

% ---- Korean fonts -----------------------------------------------------------
\input ko-font.tex      % loads luatexko, sets the hangul fonts, runs \tenpoint
\fontdimen7\tentex=0pt  % restore cmtex10's "no extra space after sentence"
                        % (ko-font.tex re-\font's \tentex, dropping gwebmac's setting)

% ---- Korean wording (the GWEB analogue of cwebman's localization list) ------
% Cross-reference notes printed beneath a code definition. Korean has no
% singular/plural inflection, so the "section"/"sections" pair reads the same.
\def\GU#1{\Gnote{이 코드는 #1절에서 쓰인다}}
\def\GUs#1{\Gnote{이 코드는 #1절에서 쓰인다}}
\def\GA#1{\Gnote{이 코드는 #1절에도 정의되어 있다}}
\def\GAs#1{\Gnote{이 코드는 #1절에도 정의되어 있다}}
% The usage note beside an entry in the section-name list.
\def\GNused#1{#1절에서 쓰임}
\def\GNuseds#1{#1절에서 쓰임}
% The running head over the list of section names, and the contents columns.
\def\Gnamesheading{절 이름 목록}
\def\Gsectionword{절}
\def\Gpageword{쪽}

% ---- LuaTeX PDF back end: blue links and Korean (UTF-16BE) bookmarks --------
% gwebmac's \Gdest/\Glink/\Gbookmark are written for pdfTeX primitives that
% luatex does not provide; we restate them with luatex's \pdfextension forms.
% A PDF outline title must be a UTF-16BE string, so Korean titles are converted
% by a small Lua helper, gweb_pdfutf16(s), which prints s as the body of a PDF
% literal string: the UTF-16BE byte-order mark followed by each code unit as a
% pair of octal-escaped bytes (surrogate-paired above the BMP).
%
% The helper is built inline. Two cautions about Lua inside \directlua: (1) the
% TeX-special characters % # _ & ^ ~ $ are made catcode 12 below (% is Lua's
% modulo operator, not a TeX comment; # is its length operator), and the
% backslash is produced as string.char(92) so none appears in the source; and
% (2) \directlua turns the source's newlines into spaces, so a Lua "--" comment
% would run to the end of the whole chunk -- hence there are none here.
\begingroup
\catcode`\%=12 \catcode`\#=12 \catcode`\_=12 \catcode`\&=12
\catcode`\^=12 \catcode`\~=12 \catcode`\$=12
\directlua{
  local function oct(v)
    local d1 = math.floor(v / 64); v = v % 64
    local d2 = math.floor(v / 8); local d3 = v % 8
    return tostring(d1) .. tostring(d2) .. tostring(d3)
  end
  function gweb_pdfutf16(s)
    local bs = string.char(92)
    local t = { bs .. oct(254) .. bs .. oct(255) }
    for _, c in utf8.codes(s) do
      if c >= 65536 then
        c = c - 65536
        local w1 = 55296 + math.floor(c / 1024)
        local w2 = 56320 + (c % 1024)
        t[#t + 1] = bs .. oct(math.floor(w1 / 256)) .. bs .. oct(w1 % 256)
        t[#t + 1] = bs .. oct(math.floor(w2 / 256)) .. bs .. oct(w2 % 256)
      else
        t[#t + 1] = bs .. oct(math.floor(c / 256)) .. bs .. oct(c % 256)
      end
    end
    tex.sprint(-2, table.concat(t))
  end
}
\endgroup
\def\Gdest#1{\pdfextension dest num #1 fith\relax}
\def\Glink#1#2{\pdfextension startlink attr{/Border[0 0 0]}goto num #1\relax
  \pdfextension literal{0 0 1 rg}#2\pdfextension literal{0 g}\pdfextension endlink}
\def\Gbookmark#1#2#3#4{\pdfextension outline goto num #2 count #3
  {\directlua{gweb_pdfutf16("\luaescapestring{#4}")}}}
% Open the bookmark (outline) pane when the document is first shown.
\pdfextension catalog{/PageMode /UseOutlines}
